%% === Tunable parameters (all diploma-specific tunables in one place) =======
% Основные параметры оформления согласно требованиям Университета ИТМО:
% - формат A4, шрифт Times New Roman, размер 14 пт, межстрочный интервал 1,5;
% - поля: левое 30 мм, правое 15 мм, верхнее 20 мм, нижнее 20 мм;
% - абзацный отступ 1,25 см.
% При необходимости поменять эти параметры изменяйте только этот файл.

%% --- Работа с русским языком ----------------------------------------------
\usepackage{cmap}                 % поиск в PDF
\usepackage{mathtext}             % русские буквы в формулах
\usepackage[T2A]{fontenc}         % кодировка шрифта
\usepackage[utf8]{inputenc}       % кодировка исходного текста
\usepackage[english,russian]{babel}

%% --- Геометрия страницы ---------------------------------------------------
\usepackage[a4paper,left=30mm,right=15mm,top=20mm,bottom=20mm,bindingoffset=0mm]{geometry}

%% --- Шрифты и размер ------------------------------------------------------
% Пакет fontspec и Times New Roman недоступны в обычной сборке pdflatex,
% поэтому применяется семейство шрифтов, совместимое с русской кодировкой T2A,
% максимально близкое к Times New Roman по начертанию.
\usepackage{mathptmx}             % аналог Times для основного текста и формул

% Базовый размер 12 pt задан в \documentclass, для основного текста задаём 14 pt.
\usepackage{anyfontsize}
\renewcommand{\normalsize}{\fontsize{14pt}{21pt}\selectfont}
\normalsize

%% --- Межстрочный интервал 1,5 ---------------------------------------------
\usepackage{setspace}
\onehalfspacing
\setlength{\emergencystretch}{3em}
\sloppy

%% --- Абзацный отступ 1,25 см, первый абзац тоже с отступом ----------------
\setlength{\parindent}{1.25cm}
\usepackage{indentfirst}

%% --- Математика ------------------------------------------------------------
\usepackage{amsmath,amsfonts,amssymb,amsthm,mathtools}
\usepackage{icomma}
\usepackage{euscript}
\usepackage{mathrsfs}

%% --- Списки ---------------------------------------------------------------
\usepackage{enumitem}
\makeatletter
\AddEnumerateCounter{\asbuk}{\russian@alph}{щ}
\makeatother
\setlist{nosep}

%% --- Рисунки --------------------------------------------------------------
\usepackage{graphicx}
\graphicspath{{../images/}{images/}}
\usepackage{caption}
\captionsetup{justification=centering,labelsep=endash,font=normalsize}
\captionsetup[figure]{name=Рисунок}
\captionsetup[table]{name=Таблица}
\usepackage{float}
\usepackage{wrapfig}
\setlength\fboxsep{3pt}
\setlength\fboxrule{1pt}

%% --- Таблицы --------------------------------------------------------------
\usepackage{array,tabularx,tabulary,booktabs}
\usepackage{longtable}
\usepackage{multirow}
\usepackage{makecell}
\usepackage{hhline}

%% --- Колонтитулы и нумерация страниц --------------------------------------
\usepackage{fancyhdr}
\pagestyle{fancy}
\fancyhf{}
\fancyfoot[C]{\thepage}
\renewcommand{\headrulewidth}{0pt}
\renewcommand{\footrulewidth}{0pt}

%% --- Плавающий знак в формулах (правило Львовского) ----------------------
\newcommand*{\hm}[1]{#1\nobreak\discretionary{}{\hbox{$\mathsurround=0pt #1$}}{}}

%% --- Дополнительное оформление --------------------------------------------
\usepackage{gensymb}
\usepackage{listings}
\lstset{
    basicstyle=\ttfamily\small,
    columns=flexible,
    breaklines=true,
    keepspaces=true
}

%% --- Гиперссылки ----------------------------------------------------------
\usepackage[unicode,pdftex,hidelinks]{hyperref}

%% --- Оформление заголовков секций (ГОСТ / ИТМО) ---------------------------
\usepackage{titlesec}
% Разделы (нумерованные главы) — прописными буквами, полужирным, по центру.
\titleformat{\section}
  {\normalfont\fontsize{14pt}{18pt}\bfseries\MakeUppercase}
  {\thesection}{1em}{}
\titlespacing*{\section}{\parindent}{2\baselineskip}{1\baselineskip}

% Подразделы — полужирным, строчными с прописной.
\titleformat{\subsection}
  {\normalfont\fontsize{14pt}{18pt}\bfseries}
  {\thesubsection}{1em}{}
\titlespacing*{\subsection}{\parindent}{1\baselineskip}{0.5\baselineskip}

% Подпункты
\titleformat{\subsubsection}
  {\normalfont\fontsize{14pt}{18pt}\bfseries}
  {\thesubsubsection}{1em}{}
\titlespacing*{\subsubsection}{\parindent}{0.75\baselineskip}{0.25\baselineskip}

% После номеров разделов/подразделов в ГОСТ 7.32 точка не ставится.
\renewcommand\thesection{\arabic{section}}
\renewcommand\thesubsection{\thesection.\arabic{subsection}}
\renewcommand\thesubsubsection{\thesubsection.\arabic{subsubsection}}

%% --- Нумерация рисунков, таблиц, формул в пределах раздела ----------------
\usepackage{chngcntr}
\counterwithin{figure}{section}
\counterwithin{table}{section}
\counterwithin{equation}{section}

%% --- Оформление неномерованных структурных элементов ----------------------
% Команда для оформления «СОДЕРЖАНИЕ», «ВВЕДЕНИЕ», «ЗАКЛЮЧЕНИЕ» и т.п.:
% центр, прописные, без нумерации, с добавлением в оглавление.
\newcommand{\structuralelement}[1]{%
    \clearpage
    \phantomsection
    \addcontentsline{toc}{section}{\MakeUppercase{#1}}%
    \begin{center}
        {\normalfont\fontsize{14pt}{18pt}\bfseries\MakeUppercase{#1}}
    \end{center}
    \vspace{0.5\baselineskip}%
}

%% --- Настройка оглавления -------------------------------------------------
\usepackage{tocloft}
\renewcommand{\cfttoctitlefont}{\hfill\normalfont\fontsize{14pt}{18pt}\bfseries\MakeUppercase}
\renewcommand{\cftaftertoctitle}{\hfill}
\renewcommand{\cftsecfont}{\normalfont\normalsize}
\renewcommand{\cftsecpagefont}{\normalfont\normalsize}
\renewcommand{\cftsubsecfont}{\normalfont\normalsize}
\renewcommand{\cftsubsecpagefont}{\normalfont\normalsize}
\renewcommand{\cftdotsep}{1.5}
\renewcommand{\cftsecleader}{\cftdotfill{\cftdotsep}}
\setlength{\cftbeforesecskip}{4pt}

%% --- Макрос для вставки рисунков с защитой от отсутствия файла ------------
\newcommand{\reportfig}[4][0.8]{%
    % #1 -- ширина (опц., доля от \textwidth), #2 -- имя файла,
    % #3 -- подпись, #4 -- метка.
    \begin{figure}[H]
        \centering
        \includegraphics[width=#1\textwidth]{#2}%
        \caption{#3}\label{#4}
    \end{figure}
}
